% ============================================================================
% sec_evaluation.tex  --  Section 3 Evaluation (Mirul / Amirul Iman), Phase 8 M8.2-M8.7
% \input under "Evaluation" in the Overleaf main file. Needs fig_dice_iou.pdf uploaded.
% Every number has its source file in a comment. Parts marked "% OPTIONAL" can be cut
% first if the report is over 4 pages.
% ============================================================================

% ---------------------------------------------------------------- M8.2
\subsection{Evaluation protocol}
All choices were made on a frozen validation split of 141 of the 707 training files (27-day time blocks, so no
validation day is also a training day). We score with the official Kaggle PQ: every annotator set of a file is its
own ground-truth entry (233 entries), a prediction matches a ground-truth filament at IoU $>0.5$, and TP, FP and FN
are pooled over all entries.
% results/phase8/Mirul/rubric_exp101.json (num_entries 233)
Our implementation gives the same value as the organisers' notebook on the same files (0.3296 both).
% SelfEvaluationResult/FINDINGS.md
Dice is \texttt{torchmetrics} \texttt{DiceScore} on the binary foreground mask (union of all filaments), averaged
over entries; the organisers' exact settings are not published.

The labels are noisy: two annotators of the same image reach only PQ 0.336 against each other, and an exact copy of
one annotator scores 0.739 on our validation entries, so the reachable PQ is far below 1.
% output/eda/eda_summary.json (0.336); SelfEvaluationResult/FINDINGS.md (0.7394)
To separate real gains from chance we used a paired bootstrap over the 141 files (2,000 draws): a gain smaller than
about 0.010 PQ is not distinguishable from noise, and retraining with another random seed changed PQ by only 0.001.
% results/phase6/Mirul/bootstrap.json; experiments/ablation_table.md Phase 6 M1a, S1
The public leaderboard (about half of the 180 test images, two decimals) has a standard deviation of about 0.014,
so we used it only as a check. Validation PQ was 0.05--0.06 above the leaderboard on every upload, but the ranking
of the steps in Table~\ref{tab:journey} was the same on both.
% experiments/ablation_table.md Phase 6 M1a (LB SD 0.013-0.015); experiments/log.csv (lb_score column)

% ---------------------------------------------------------------- M8.3 + M8.1
\subsection{Results}
Table~\ref{tab:journey} shows how each step changed the score. The largest gains came from native 2048-pixel input
($+0.025$), from combining two models ($+0.010$) and from the four flipped views with two-stage fusion ($+0.011$).
The last two are above the noise level: $+0.039$ [0.027, 0.050] for the ensemble against the 1024-pixel model and
$+0.012$ [0.004, 0.020] for the final pipeline against the plain ensemble (95\% bootstrap intervals).
% results/phase6/Mirul/bootstrap.json (exp010 vs exp006); results/phase6/Mirul/bootstrap_m2.json (exp101 vs exp010)

\begin{table}[t]
\centering
\caption{From baseline to final pipeline. Validation: official PQ and mean Dice on the frozen 141-file split.
LB: public leaderboard, about 50\% of the test set.}
\label{tab:journey}
\small
\begin{tabular}{@{}lccc@{}}
\hline
Step & Val PQ & Dice & LB \\
\hline
Mask R-CNN R50, 1024 px (baseline)     & 0.391 & 0.609 & 0.34 \\ % log.csv exp003
+ all annotator sets as training data  & 0.398 & 0.643 & 0.35 \\ % log.csv exp006
+ native 2048 px, R50                  & 0.423 & 0.648 & 0.36 \\ % log.csv exp007 (thr 0.6)
\quad X101 instead of R50              & 0.426 & 0.642 & 0.37 \\ % log.csv exp008
+ ensemble X101 + R50                  & 0.437 & 0.651 & 0.38 \\ % log.csv exp010
+ 56$\times$56 mask head, 4 flip views, & \textbf{0.448} & \textbf{0.652} & \textbf{0.39} \\ % log.csv exp101
\quad two-stage fusion (\textbf{final}) & & & \\
\hline
\end{tabular}
\end{table}

\textbf{Final pipeline.} On validation it reaches PQ $0.448$ (95\% interval $0.425$--$0.470$), with segmentation
quality SQ $0.673$ and recognition quality RQ $0.665$ (1,083 TP, 495 FP, 595 FN), and mean Dice $0.652$
($0.628$--$0.672$). The intervals come from 2,000 bootstrap draws of the 141 files.
% results/phase8/Mirul/rubric_exp101.json (PQ, SQ, RQ, TP/FP/FN, Dice_samplewise, bootstrap_95ci_files)
Fig.~\ref{fig:diceiou} shows the Dice and IoU of every pair of a predicted and a ground-truth filament that overlap.
The matched pairs are tightly grouped (median Dice 0.81, median IoU 0.67); the long tail below the 0.5 cut-off are
filaments we found but outlined too differently from the annotator.
% rubric_exp101.json dist_Dice_matched_TP.median 0.806, dist_IoU_matched_TP.median 0.675
\textbf{One-to-many and many-to-one.} Of the ground-truth filaments, 20 were split into several predicted pieces
(19 into two, one into three); 58 of our 933 masks cover two or more ground-truth filaments (55 cover two, 3 cover
three). A piece counts when at least half of its area lies inside the other mask. Merging is therefore about three
times more common than fragmentation.
% rubric_exp101.json one_to_many_*, many_to_one_*, n_pred_masks 933

\begin{figure}[t]
\centering
\includegraphics[width=\columnwidth]{fig_dice_iou.pdf}
\caption{Dice (left) and IoU (right) of every overlapping pair of predicted and ground-truth filaments on
validation (all annotator sets). Dark: pairs counted as true positives (IoU $>0.5$, dashed line).}
\label{fig:diceiou}
% docs/report/m81_rubric.py -> docs/report/figures/fig_dice_iou.pdf
\end{figure}

% ---------------------------------------------------------------- M8.4
\subsection{What did not work}
Table~\ref{tab:negative} lists ideas that did not pass the 0.010 noise bar. Training-recipe changes were tested on
the cheaper 1024-pixel R50 model (reference PQ 0.398), all others on the 2048-pixel models.

\begin{table}[t]
\centering
\caption{Rejected ideas (validation official PQ; $\Delta$ against the reference of that test).}
\label{tab:negative}
\small
\begin{tabular}{@{}lcl@{}}
\hline
Idea & $\Delta$PQ & Reason \\
\hline
CLAHE + solar-disk mask      & $+0.000$ & no gain \\                    % ablation_table Retraining A
SAM 2 mask refinement        & $-0.019$ & tighter than annotators \\    % ablation_table Phase 5 step 1
Mask2Former (Swin-S)         & $-0.024$ & twice the merges \\           % ablation_table Phase 5 step 2
3$\times$ longer schedule    & $+0.001$ & noise, lower Dice \\          % ablation_table S2 (exp202)
90$^\circ$ rotations + zoom  & $-0.001$ & no gain \\                    % ablation_table S3 (exp203)
Copy-paste of filaments      & $-0.004$ & more false alarms \\          % ablation_table S4 (exp204)
Self-training, 2,098 images  & $-0.002$ & model too cautious \\         % ablation_table S5 (exp205)
Third ensemble member        & $-0.003$ & same errors as X101 \\        % ablation_table Phase 7 prep A
Learned re-scoring           & $-0.006$ & too few labels \\             % ablation_table Phase 7 prep D (exp016)
Retrain on all 707 files     & LB $=$   & 95\% same masks \\            % ablation_table Phase 7 prep B (exp015)
\hline
\end{tabular}
\end{table}

% OPTIONAL (cut first): one-sentence explanations of Table 2
SAM 2 draws the dark core of a filament, about 15\% smaller than the human outline, which confirms that the labels,
not mask sharpness, limit SQ. Self-training used only confident pseudo-labels (score $\geq 0.85$), so faint
filaments in the extra images were taught as background.
% ablation_table Phase 5 step 1, Phase 6 S5

% ---------------------------------------------------------------- M8.5
\subsection{Error analysis}
We put every validation error of the final pipeline into one cause and computed the PQ if only that cause were
fixed. The largest is confidence: 282 ground-truth filaments are found by a raw prediction (IoU $>0.5$) but scored
below the 0.6 cut-off ($+0.072$ PQ if kept; 233 of them medium-sized). Next are 305 false alarms that no annotator
drew ($+0.046$), 162 filaments with no prediction near them ($+0.043$) and 115 near misses with IoU 0.3--0.5
($+0.031$). A further 190 ``false alarms'' were drawn by another annotator of the same image, i.e.\ label noise
rather than model error.
% results/phase6/Mirul/error_budget_exp101.json; experiments/ablation_table.md Phase 6 M4
By size, medium filaments score PQ 0.430 against 0.493 for large ones; the 12 small filaments ($<200$\,px$^2$)
are all missed, which costs at most 0.005 PQ. By station, PQ ranges from 0.418 (El Teide, most false alarms) to
0.468 (Big Bear); with only 20--28 files per station these gaps are within noise.
% error_budget_exp101.json by size / station; docs/HANDOVER.md section 6 (small <= 0.005)

% ---------------------------------------------------------------- M8.6
\subsection{Runtime}
Inference runs eight forward passes per image (two models $\times$ four views) plus fusion and post-processing:
median 6.0\,s and 95th percentile 8.2\,s per 2048$\times$2048 image on one NVIDIA L4 GPU, 18 minutes for all 180
test images, with 2.9\,GB peak GPU memory. Training at 2048 pixels needs about 20\,GB and was done on an A100.
% results/phase7/final_run/report_test.json; experiments/log.csv exp008 (peak 20.2 GB)

% ---------------------------------------------------------------- M8.7
\subsection{Strengths, weaknesses and improvements}
\textbf{Strengths.} Every step was chosen on a fixed split against a measured noise level, so the gains carried
over to the leaderboard (0.34 to 0.39). The method is simple (two standard Mask R-CNN models, no pre-processing) and
a clean clone of the repository reproduces the validation result exactly.
% results/phase7/final_run/report_val.json (PQ 0.4480 again)
\textbf{Weaknesses.} The main loss is confidence near the cut-off: many correct medium-sized filaments are scored
just below 0.6, and one threshold cannot keep them without adding false alarms. Over-merging of nearby filaments is
three times more common than fragmentation, and small filaments are not found.
\textbf{Improvements.} Better calibrated scores for medium filaments, for example from more labelled data or
self-training with a lower pseudo-label threshold, which we did not have time to test; and a rule or model to
split masks that join two filaments.
% error_budget_exp101.json (fn_low_score); ablation_table Phase 6 S5 (lower threshold not tried)
